\ifajar
  \section{Latihan}
\else
  \section{Bab 11: Arsitektur Frontend Berskala Besar}
\fi

Seluruh file yang dipakai pada latihan ini tersedia di
\berkas{sources/chapter11/}, dihitung relatif terhadap akar repositori.
Latihan ini membutuhkan Python 3 dan satu browser berbasis Chromium, tanpa
Docker dan tanpa pustaka tambahan.

Seluruh permintaan diarahkan ke aplikasi yang berjalan di komputer sendiri,
dan sifatnya mengamati, bukan membebani. Angka milidetik tiap mahasiswa akan
berbeda, karena bergantung pada komputernya. Yang dilaporkan adalah pola
perbandingannya, bukan angka persisnya.

Penyiapannya satu perintah, seperti pada Listing \ref{lst:c11-siapkan}.

\begin{lstlisting}[language=bash, caption={Menyalakan aplikasi katalog, tersedia di \berkas{sources/chapter11/aplikasi.py}}, label={lst:c11-siapkan}]
cd sources/chapter11
python3 aplikasi.py &

# Keluarannya:
# Melayani 200 buku di http://127.0.0.1:8080/penuh  (Ctrl+C untuk berhenti)
\end{lstlisting}

\ifajar
  \subsection*{Latihan 1: Menghitung Komponen dan Biaya State}
\else
  \subsection*{Prosedur 1: Menghitung Komponen dan Biaya State}
\fi

Latihan ini membuktikan bahwa jumlah pulau tidak tumbuh mengikuti besar
halaman, sedangkan jumlah komponen dan ukuran halaman tumbuh.

\begin{enumerate}
  \item Nyalakan aplikasi dengan jumlah buku bawaan, lalu hitung kemunculan
        \texttt{data-komponen} dan \texttt{data-pulau} pada kedua varian,
        seperti pada Listing \ref{lst:c11-lat1-jalan}.
  \item Catat ukuran HTML kedua varian, lalu hitung selisihnya.
  \item Hentikan aplikasi, nyalakan ulang dengan
        \texttt{JUMLAH\_BUKU=1000}, lalu ulangi kedua langkah di atas.
\ifajar
  \item Isi \tabisi{tab:c11-lat1-hasil}. Tabel \ref{tab:c11-lat1-contoh}
        memperlihatkan contoh yang sudah terisi.
\else
  \item Bandingkan hasil pengukuran sendiri dengan Tabel
        \ref{tab:c11-lat1-contoh}.
\fi
  \item Jawab dua pertanyaan berikut. Pertama, berapa persen tambahan ukuran
        HTML pada varian penuh, dan dari mana asalnya? Kedua, mengapa jumlah
        pulau tetap dua walaupun bukunya lima kali lebih banyak?
\end{enumerate}

\begin{lstlisting}[language=bash, caption={Menghitung komponen dan ukuran halaman, tersedia di \berkas{sources/chapter11/aplikasi.py}}, label={lst:c11-lat1-jalan}]
curl -s http://127.0.0.1:8080/penuh | grep -c 'data-komponen='
curl -s http://127.0.0.1:8080/selektif | grep -c 'data-pulau='
curl -s http://127.0.0.1:8080/penuh | wc -c
curl -s http://127.0.0.1:8080/selektif | wc -c

# Keluarannya:
# 403
# 2
# 71950
# 55674
\end{lstlisting}

\begin{table}[!htb]
  \centering
  \caption{Contoh hasil Latihan 1 yang sudah terisi}
  \label{tab:c11-lat1-contoh}
  \begin{tabularx}{\textwidth}{Xllll}
    \hline
    \textbf{Buku} & \textbf{Komponen} & \textbf{Pulau} & \textbf{HTML penuh} & \textbf{HTML selektif} \\
    \hline
    200 & 403 & 2 & 71.950 & 55.674 \\
    \hline
    1.000 & 2.003 & 2 & 358.360 & 275.679 \\
    \hline
  \end{tabularx}
\end{table}

\ifajar
\begin{table}[!htb]
  \centering
  \caption{Lembar isian hasil pengukuran Latihan 1}
  \label{tab:c11-lat1-hasil}
  \begin{tabularx}{\textwidth}{Xllll}
    \hline
    \textbf{Buku} & \textbf{Komponen} & \textbf{Pulau} & \textbf{HTML penuh} & \textbf{HTML selektif} \\
    \hline
    200 & \dots & \dots & \dots & \dots \\
    \hline
    1.000 & \dots & \dots & \dots & \dots \\
    \hline
  \end{tabularx}
\end{table}
\fi

\ifajar
  \subsection*{Latihan 2: Biaya Hidrasi Dua Varian}
\else
  \subsection*{Prosedur 2: Biaya Hidrasi Dua Varian}
\fi

Latihan ini mengukur biaya hidrasi kedua varian pada tiga besar halaman, lalu
memisahkan bagian yang tumbuh dari bagian yang tetap.

\begin{enumerate}
  \item Jalankan \berkas{ukur-hidrasi.py} pada jumlah buku bawaan, seperti
        pada Listing \ref{lst:c11-lat2-jalan}.
  \item Hentikan aplikasi, nyalakan ulang dengan
        \texttt{JUMLAH\_BUKU=1000}, lalu ukur sekali lagi. Ulangi untuk
        \texttt{JUMLAH\_BUKU=2000}.
\ifajar
  \item Isi \tabisi{tab:c11-lat2-hasil}. Tabel \ref{tab:c11-lat2-contoh}
        memperlihatkan contoh yang sudah terisi.
\else
  \item Bandingkan hasil pengukuran sendiri dengan Tabel
        \ref{tab:c11-lat2-contoh}.
\fi
  \item Buka kedua varian di browser, yaitu \texttt{/penuh} dan
        \texttt{/selektif} pada port 8080, lalu bandingkan baris laporan di
        bawah judul halamannya.
  \item Jawab tiga pertanyaan berikut. Pertama, bagian mana yang tumbuh
        mengikuti besar halaman, dan bagian mana yang tidak? Kedua, mengapa
        angka varian selektif tidak urut menurut besar halaman? Ketiga, pada
        besar halaman berapa selisihnya mulai pantas dipersoalkan, dan dengan
        alasan apa?
\end{enumerate}

\begin{lstlisting}[language=bash, caption={Mengukur biaya hidrasi lima kali tiap varian, tersedia di \berkas{sources/chapter11/ukur-hidrasi.py}}, label={lst:c11-lat2-jalan}]
python3 ukur-hidrasi.py

# Keluarannya:
#   Varian     Buku  Komponen  HTML byte  JS byte  Tengah ms   Rentang ms
#   penuh       200       403      71950     3723       1.50    1.20-2.10
#   selektif    200         2      55674     3278       0.50    0.50-0.70
\end{lstlisting}

\begin{table}[!htb]
  \centering
  \caption{Contoh hasil Latihan 2 yang sudah terisi}
  \label{tab:c11-lat2-contoh}
  \begin{tabularx}{\textwidth}{Xllll}
    \hline
    \textbf{Buku} & \textbf{Penuh, ms} & \textbf{Selektif, ms} & \textbf{Komponen penuh} & \textbf{Pulau} \\
    \hline
    200 & 1,50 & 0,50 & 403 & 2 \\
    \hline
    1.000 & 3,90 & 1,00 & 2.003 & 2 \\
    \hline
    2.000 & 7,10 & 0,80 & 4.003 & 2 \\
    \hline
  \end{tabularx}
\end{table}

\ifajar
\begin{table}[!htb]
  \centering
  \caption{Lembar isian hasil pengukuran Latihan 2}
  \label{tab:c11-lat2-hasil}
  \begin{tabularx}{\textwidth}{Xllll}
    \hline
    \textbf{Buku} & \textbf{Penuh, ms} & \textbf{Selektif, ms} & \textbf{Komponen penuh} & \textbf{Pulau} \\
    \hline
    200 & \dots & \dots & \dots & \dots \\
    \hline
    1.000 & \dots & \dots & \dots & \dots \\
    \hline
    2.000 & \dots & \dots & \dots & \dots \\
    \hline
  \end{tabularx}
\end{table}
\fi

\ifajar
  \subsection*{Latihan 3: Biaya Perpindahan pada Kerangka}
\else
  \subsection*{Prosedur 3: Biaya Perpindahan pada Kerangka}
\fi

Latihan ini membandingkan biaya satu perpindahan lewat fragmen dengan biaya
memuat dokumen penuh, lalu memeriksa perilaku kerangkanya di browser.

\begin{enumerate}
  \item Jalankan \berkas{ukur-rute.py}, seperti pada Listing
        \ref{lst:c11-lat3-jalan}, lalu catat keempat barisnya.
  \item Buka \texttt{http://127.0.0.1:8080/kerangka/katalog} di browser,
        tekan tautan "Buku 3", lalu perhatikan alamat di bilah alamat dan
        bagian halaman yang berubah.
  \item Tekan tombol kembali browser, lalu catat apakah katalognya pulih.
  \item Tekan tombol Tambah dua kali, lalu pindah ke halaman buku dan
        kembali lagi. Catat angka keranjangnya sesudah kembali.
\ifajar
  \item Isi \tabisi{tab:c11-lat3-hasil}. Tabel \ref{tab:c11-lat3-contoh}
        memperlihatkan contoh yang sudah terisi.
\else
  \item Bandingkan hasil pengukuran sendiri dengan Tabel
        \ref{tab:c11-lat3-contoh}.
\fi
  \item Jawab tiga pertanyaan berikut. Pertama, mengapa fragmen katalog
        hampir sebesar dokumennya? Kedua, apa yang tetap dihemat perpindahan
        lewat fragmen walaupun selisih bytenya kecil? Ketiga, mengapa angka
        keranjang kembali ke nol, dan di mana seharusnya \textit{state} itu
        disimpan?
\end{enumerate}

\begin{lstlisting}[language=bash, caption={Membandingkan ukuran dokumen dan fragmen, tersedia di \berkas{sources/chapter11/ukur-rute.py}}, label={lst:c11-lat3-jalan}]
python3 ukur-rute.py

# Keluarannya:
#   Tujuan                    Status   Dokumen   Fragmen  Fragmen %
#   /katalog                 200/200     55943     55089      98.5%
#   /katalog?saring=jaringan 200/200     10414      9560      91.8%
#   /buku/3                  200/200      1071       217      20.3%
#   /gagal                   500/500      1023       169      16.5%
\end{lstlisting}

\begin{table}[!htb]
  \centering
  \caption{Contoh hasil Latihan 3 yang sudah terisi}
  \label{tab:c11-lat3-contoh}
  \begin{tabularx}{\textwidth}{Xll}
    \hline
    \textbf{Hal yang diamati} & \textbf{Hasil} & \textbf{Keterangan} \\
    \hline
    Fragmen \texttt{/buku/3} & 217 byte & dokumennya 1.071 byte \\
    \hline
    Fragmen \texttt{/katalog} & 55.089 byte & 98,5 persen dari dokumennya \\
    \hline
    Alamat sesudah menekan tautan & \texttt{/kerangka/buku/3} & tanpa muat ulang \\
    \hline
    Tombol kembali browser & katalog pulih & 200 baris \\
    \hline
    Keranjang sesudah kembali & 0 & disimpan di DOM wilayah \\
    \hline
  \end{tabularx}
\end{table}

\ifajar
\begin{table}[!htb]
  \centering
  \caption{Lembar isian hasil pengamatan Latihan 3}
  \label{tab:c11-lat3-hasil}
  \begin{tabularx}{\textwidth}{Xll}
    \hline
    \textbf{Hal yang diamati} & \textbf{Hasil} & \textbf{Keterangan} \\
    \hline
    Fragmen \texttt{/buku/3} & \dots & \dots \\
    \hline
    Fragmen \texttt{/katalog} & \dots & \dots \\
    \hline
    Alamat sesudah menekan tautan & \dots & \dots \\
    \hline
    Tombol kembali browser & \dots & \dots \\
    \hline
    Keranjang sesudah kembali & \dots & \dots \\
    \hline
  \end{tabularx}
\end{table}
\fi

\ifajar
  \subsection*{Latihan 4: Apa yang Tersisa Tanpa Skrip}
\else
  \subsection*{Prosedur 4: Apa yang Tersisa Tanpa Skrip}
\fi

Latihan ini membuktikan bahwa isi katalog dan penyaringnya tidak bergantung
pada JavaScript, dan memperlihatkan batas \textit{error} menahan satu wilayah
yang gagal.

\begin{enumerate}
  \item Jalankan \berkas{ukur-tanpa-js.py}, seperti pada Listing
        \ref{lst:c11-lat4-jalan}, lalu catat keempat barisnya.
  \item Buka \texttt{http://127.0.0.1:8080/kerangka/katalog?blokir\_skrip=1}
        di browser, ketik "jaringan" di kotak penyaring, lalu tekan tombol
        Saring. Catat jumlah baris yang tersisa dan apakah halamannya dimuat
        ulang.
  \item Buka alamat yang sama tanpa \texttt{blokir\_skrip}, lalu ulangi
        penyaringan yang sama. Catat perbedaannya.
  \item Tekan tautan "Bagian yang gagal", lalu catat apa yang tampil di
        wilayah dan apakah navigasinya masih utuh.
\ifajar
  \item Isi \tabisi{tab:c11-lat4-hasil}. Tabel \ref{tab:c11-lat4-contoh}
        memperlihatkan contoh yang sudah terisi.
\else
  \item Bandingkan hasil pengukuran sendiri dengan Tabel
        \ref{tab:c11-lat4-contoh}.
\fi
  \item Jawab tiga pertanyaan berikut. Pertama, bagian mana yang hilang
        tanpa skrip, dan bagian mana yang tetap? Kedua, apa yang akan terlihat
        di halaman ini bila katalognya di-\textit{render} di browser seperti varian SPA
        (Single-Page Application) pada Bab 8? Ketiga, mengapa router
        memeriksa status responsnya, bukan cuma menangkap kegagalan jaringan?
\end{enumerate}

\begin{lstlisting}[language=bash, caption={Memeriksa isi yang tetap tampil tanpa skrip, tersedia di \berkas{sources/chapter11/ukur-tanpa-js.py}}, label={lst:c11-lat4-jalan}]
python3 ukur-tanpa-js.py

# Keluarannya:
#   Halaman                            Skrip jalan Skrip mati  Form
#   /selektif                                  200        200   ada
#   /selektif?saring=jaringan                   33         33   ada
#   /kerangka/katalog                          200        200   ada
#   /kerangka/katalog?saring=jaringan           33         33   ada
\end{lstlisting}

\begin{table}[!htb]
  \centering
  \caption{Contoh hasil Latihan 4 yang sudah terisi}
  \label{tab:c11-lat4-contoh}
  \begin{tabularx}{\textwidth}{Xll}
    \hline
    \textbf{Hal yang diamati} & \textbf{Skrip jalan} & \textbf{Skrip mati} \\
    \hline
    Baris katalog & 200 & 200 \\
    \hline
    Baris sesudah disaring & 33 & 33 \\
    \hline
    Cara penyaringan & tanpa muat ulang & satu permintaan baru \\
    \hline
    Wilayah yang gagal & diganti pesan & halaman 500 dari server \\
    \hline
    Navigasi sesudah gagal & utuh, 4 tautan & utuh, 4 tautan \\
    \hline
  \end{tabularx}
\end{table}

\ifajar
\begin{table}[!htb]
  \centering
  \caption{Lembar isian hasil pengamatan Latihan 4}
  \label{tab:c11-lat4-hasil}
  \begin{tabularx}{\textwidth}{Xll}
    \hline
    \textbf{Hal yang diamati} & \textbf{Skrip jalan} & \textbf{Skrip mati} \\
    \hline
    Baris katalog & \dots & \dots \\
    \hline
    Baris sesudah disaring & \dots & \dots \\
    \hline
    Cara penyaringan & \dots & \dots \\
    \hline
    Wilayah yang gagal & \dots & \dots \\
    \hline
    Navigasi sesudah gagal & \dots & \dots \\
    \hline
  \end{tabularx}
\end{table}
\fi
